\section{Stencil Computations}

\subsection{Introduction}

A \textit{stencil} is a function used to compute the value of a point in a structured grid for some timestep by utilizing values from nearby points at prior timesteps \cite{Frigo2005Stencil}.
A \textit{ stencil computation} applies a given stencil to all points of a grid for a given number of timesteps.

Stencil computation is a widely used paradigm, particularly in high performance computing (HPC). 
Also known as \textit{structured grid} computation, this paradigm was identified in  ``The Landscape of Parallel Computing Research: A View from Berkeley'' (2006) 
\cite{Asanovic2006Landscape}, as one of  
seven key numerical methods 
that dominate the HPC landscape. 
Its also worth noting that spectral methods 
are another of identified paradigms, 
but that will be relevant in the next section.

Stencils funtions derived from finite difference methods \cite{Taflove2005Ch9, }. 

Others like lattice boltzmann, or cellular automata.

Fields that use stencils.




Explain the premise of stencil computations, and where they are found.

Explain our formalism of $N$ cells, $T$ time steps.

Applications

Geology \cite{Dursun2009Stencil}

Oceanography \cite{Rainer1992Atlantic}

Fluid Dynamics \cite{Blazek2015Ch1, Ferziger2020, Kruger2017LBM}.

Electromagnetics \cite{Taflove2005Ch9}

Mechanical engineering \cite{Rappaz2003, Stathopoulos2000Material}

Meteorology \cite{Singh2023SPARTA, Avissar1989Weather, Kalnay1990Weather}

%Image Processing \cite{}
Complexity of implementing stencil computations 
has led to a proliferation of specialized stencil compilers 
\cite{Bauer2019PhaseField, Bauer2021lbmpy, Ahmad2022, Tang2011Pochoir, Holewinski2012Stencil}.

\subsubsection{Lattice Boltmann Method}

The

Aeroacoustics \cite{Casalino2014LBM}

\subsubsection{Cellular Automata}

Cellular Automata (CA) are a subset of stencil computations
with applicability to modelling natural and self-replicating systems \cite{Mendicino2015Eco, Caux2010HashLife, Sitko2016Cellular, Wolfram1984Cellular, Kyparissas2020Cellular}.
Cell values in a CA are often called states 
and typically have a small finite set of spatial values.  
As a result, different CA can generate varying degrees of regularity 
across different spatial and temporal scales.
This allows for powerful set of caching based algorithms, 
commonly known as \textit{Hash-Life}  \cite{Gosper1984HashLife}



\subsection{Direct Methods}

\subsubsection{Py Stencil LBM}
\begin{outline}
\1 $O(NT)$ work./
\1 Pystencil and LBM methods
\2 Introduced in \cite{Bauer2019PhaseField}
\2 LBMpy introduced to \cite{Bauer2021lbmpy}
\2 Expanded to CMMRT, interesting domain specific optimizations like zero centered storage \cite{Hennig2023CMMRT}
\2 Deploys to walberla \cite{Bauer2021walberla}
\2 Algebraic compilation for complex non-linear stencils applied to high dimensional grid data.
\end{outline}

\subsubsection{Cache Friendly}
\begin{outline}
  \1 Study on stencils and memory subsystems \cite{Datta2009Stencil}
  \1 Implicit and Explicit Cache optimizations for Cell processor \cite{Kamil2006Stencil} 
  \1 Another Cell processor paper \cite{DeLaCruz2010Stencil}.
  \1 Impact of modern memory sub-systems on stencils \cite{Kamil2005Stencil}
  \1 Cache blocking, prefetch, and SIMD \cite{Dursun2009Core}.

\1 Cache Oblivious methods
\2 Trapezoidal Decompositions in Pochoir, cache oblivious \cite{Frigo2005Stencil, Tang2011Pochoir}
\1 Pluto \cite{Bondhugula2008Automatic, Bondhugula2008Practical}
\1 More general methods like polyhedral compilers. 
\1 Automatic stencil tiling \cite{Li2004Tiling}
\1 Load balancing / parallelism for tiling \cite{Krishnamoorthy2007Tiling}
\end{outline}

\subsection{Approximate Methods}
Krylov methods are used to solve linear systems $Ax =b$, which is a distinct paradigm from stencil computation, though many stencil computations can be expressed in terms of a sparse linear system solve.
These methods are iterative and approximate, which the core idea being repeated matrix-vector multiplication operations \cite{Ipsen1998Krylov, Saad1989Krylov}. 
These methods are often paired with preconditions, and discrete Fourier Transforms have been applied in this context \cite{Chan1984fd, Erlangga2006, Gmeiner2013}.
These methods often require manual convergence analysis
\cite{Brown1994Krylov, Kuijlaars2006Krylov, Notay2008Krylov}.
As a result, Krylov methods are often specialized and tailored to individual problems \cite{Asgharzadeh2017Krylov, Ashrafizadeh2015, Garrappa2015Krylov,Mang2015Krylov}.

\subsection{Optimizations}

Useful paper with lots of refs on SIMD optimizations \cite{Henretty2011Stencil} 

\subsection{FPGA}
\begin{outline}
  \1 Cellular Automata \cite{Kyparissas2020Cellular}
  \1 TPU computations using convolutions with ConvStencil \cite{Chen2024ConvStencil}
  \1 SPARTA for accelerating 2D diffusion stencils on FPGAs using MLIR for weather \cite{Singh2023SPARTA}
  \1 Efficient GPU implementation model \cite{Micikevicius2009FDGPU}
  \1 Tiling / Ghost zone on GPU \cite{Meng2009Stencil}
  \1 Code generation for GPU \cite{Holewinski2012Stencil}
  \1 Multilevel parallelism for cluster computing \cite{Dursun2009Stencil}
  \1 Another cluster / multi-level optimization paper \cite{Augustin2009Stencil, Treibig2011Stencil}
  \1 Multicore optimization and tuning \cite{Datta2008Stencil}
  
\end{outline}


\subsection{Indirect Methods}

\subsubsection{Regularity Caching}
Cellular automata typically have a small number of possible states per cell.
Exploit regularities in cellular spaces.
Developed in 2D \cite{Gosper1984HashLife}.
Applied to biology and 3D \cite{Caux2010HashLife}.
Allows for the exploration of enormous spatial and time scales that couldn't be supported with explicit direct solvers, which is particularly relavent for the study of self-replicating systems \cite{Hutton2010Codd, Stevens2011Cellular}.

\subsubsection{Spectral}
\begin{outline}
  \1 Convolution based methods  \cite{Januario2016}
  \1 Gaussian based methods \cite{Ahmad2025Gaussian}
  \1 FFT Based methods \cite{Ahmad2021, Ahmad2023},
  \2 Code generation with Fourst \cite{Ahmad2022},
  \2 Variation of the aperiodic algorithm was explored with TurboStencil \cite{Liu2023Turbo}
  \2 Extended to GPU / TCU \cite{Han2025FlashFFT},
  \1 Applied to specilized non-linear problem \cite{Ahmad2024Option} 
  \1 And of course we added inhomogenous domains and time varying stencils \cite{Bentley2025Fast}
\end{outline}

